\section{Independent computation formulas}\label{tsc:computation}
\subsection{Direct three-digamma finite parts}
Suppose $0\le a_1<a_2<a_3<1$, put $a_4=a_1+1$, and let
$L_j=a_{j+1}-a_j$. On the interval beginning at $a_j$, use the coordinate
$t=x-a_j$ and offsets $d_{j\ell}=\{a_j-a_\ell\}$, with $d_{jj}=0$.
Set
\[
 c_j=\prod_{\ell\ne j}\gamma_0(d_{j\ell}).
\]
Then a directly integrable version of the canonical triple is
\begin{equation}\label{tsc:direct-triple}
 \calJ_{000}(\mathbf a)=\sum_{j=1}^3\left\{
   \int_0^{L_j}\left[\prod_{\ell=1}^3\gamma_0(t+d_{j\ell})
                              -\frac{c_j}{t}\right]\dd t
                          +c_j\log L_j\right\}.
\end{equation}
No fractional part is applied inside a single interval: all the offsets
stay in the appropriate left-limit branches until its upper endpoint.
The removable value of the bracket at zero is
\[
 \gamma c_j+
  \sum_{\ell\ne j}\gamma_0'(d_{j\ell})
             \prod_{h\ne j,\ell}\gamma_0(d_{jh}).
\]
The three-digamma product is $-\calJ_{000}$. This quadrature uses no
Tornheim function, Mellin continuation, or polylogarithm derivative, making
it independent of the right side of \eqref{tsc:psi-integrals}.

\subsection{Stable evaluation of the order-zero polylogarithm jet}
For $\mu=-x+\ii\theta$ with $|\mu|<2\pi$, differentiating the classical
Jonqui\`ere expansion \cite[Eq.~25.12.12]{tsc:dlmf-poly} gives
\begin{equation}\label{tsc:Li-order-zero}
 \Li_0^{[1]}(e^\mu)
   =-\frac{\gamma+\Log(-\mu)}{\mu}
       +\sum_{k\ge0}\frac{\zeta^{[1]}(-k)}{k!}\mu^k.
\end{equation}
The principal branch is used for $\Log(-\mu)$.
For $0\le x\le1$ and $0<|\theta|\le\pi$, this convergent disk contains
the entire evaluation region. The coefficients can be generated without
spectral differencing of a polylogarithm:
\begin{align}\label{tsc:negative-zeta-derivative}
 \zeta^{[1]}(0)&=-\frac\ell2,\notag\\
 \frac{\zeta^{[1]}(-k)}{k!}
  &=\frac{(-1)^{k/2}\zeta(k+1)}{2(2\pi)^k},
                                      &&k\ge2\text{ even},\notag\\
 \frac{\zeta^{[1]}(-k)}{k!}
  &=\frac{\zeta(-k)}{k!}
       \left(\ell-\psi(k+1)
                     -\frac{\zeta^{[1]}(k+1)}{\zeta(k+1)}\right),
                                      &&k\ge1\text{ odd}.
\end{align}
These follow directly by differentiating the Riemann functional equation
at its trivial zeros and at its negative odd arguments.
For $x>1$, use instead the rapidly convergent defining series
\begin{equation}\label{tsc:Li-log-series}
 \Li_0^{[1]}(ze^{-x})=-\sum_{n\ge1}(ze^{-x})^n\log n.
\end{equation}
The checker precomputes the first series coefficients, switches to the
second series after $x=1$, and inserts the analytic first derivative at
the removable endpoint of the subtracted Mellin integrand. The degree and
precision choices are numerical safeguards; no interval enclosure is
claimed. Six root-of-unity endpoint evaluations are also checked against
the independent finite Gamma formula \eqref{tsc:W-boundary}.

\subsection{Explicit singularity subtraction for the cubic example}
Let $f(x)$ be the integrand in \eqref{tsc:explicit-cubic-eq}, and put
\[
 r_1=-\frac{\psi(1/4)}\pi,\qquad r_2=\frac{\psi(1/2)}\pi,
 \qquad g(x)=f(x)-\frac{r_1}{x-1/4}-\frac{r_2}{x-1/2}.
\]
Then
\begin{equation}\label{tsc:cubic-numeric}
 \PV\int_0^1 f(x)\dd x=\int_0^1g(x)\dd x+r_1\log3.
\end{equation}
The removable interior values used in the implementation are
\begin{align*}
 g(1/4)&=-\frac{\psi'(1/4)}\pi-2\psi(1/4)
                         -\frac{r_2}{1/4-1/2},\\
 g(1/2)&=\frac{\psi'(1/2)}\pi-2\psi(1/2)
                         -\frac{r_1}{1/2-1/4},\\
 g(0)&=-\pi+4r_1+2r_2.
\end{align*}
Thus the numerical check uses an ordinary integral with explicitly removed
poles, rather than an extrapolated cutoff sequence.

\subsection{Finite tails and an implementation safeguard}
For the positive-integer Tornheim diagnostics, the Mellin integral is
truncated at a finite $R$. For $r,s\ge0$ integral, unit colors, and
$x\ge\log2$,
\[
 |\Li_r(ze^{-x})\Li_s(we^{-x})|\le4e^{-2x}.
\]
Consequently the omitted tail of \eqref{tsc:T-mellin}, for integer $t\ge1$,
is at most $4\Gamma(t,2R)/(2^t\Gamma(t))$. This finite-tail implementation
also avoids enormous-exponent complex logarithms that can arise when a
black-box infinite-range quadrature samples astronomically large $x$.
It is a safeguard in the newly written numerical checker, not a correction
to an existing repository theorem.

\subsection{Artifact map}
The principal reproducibility files are
\begin{center}
\begin{tabular}{p{0.43\textwidth}p{0.49\textwidth}}
\toprule
File & Purpose\\\midrule
\src{article.tex}, \src{sections/} & Modular complete proof text.\\
\src{article_standalone.tex} & Flattened single-file LaTeX source.\\
\src{article.pdf} & Compiled article.\\
\src{code/verify_exact.py} & 621 exact assertions and coefficient-table generation.\\
\src{code/calculus.py} & Subtracted quadratures and stable polylog jet utilities.\\
\src{code/verify_numerical.py} & 56 numerical comparisons with explicit tolerances.\\
\src{results/} & Executed JSON summaries and finite coefficient tables.\\
\src{notes/} & Source snapshot, audit, proof dependencies, and integration note.\\
\src{verification/} & Build and executed check logs.\\
\src{SHA256SUMS} & Integrity manifest; not a mathematical certificate.\\\bottomrule
\end{tabular}
\end{center}
